\documentclass[11pt,a4paper]{article}
\usepackage[T1]{fontenc}
\usepackage[utf8]{inputenc}
\usepackage[margin=2.4cm]{geometry}
\usepackage{booktabs}
\usepackage{array}
\usepackage{xcolor}
\usepackage{titlesec}
\usepackage{enumitem}
\usepackage{listings}
\usepackage[hidelinks]{hyperref}

\newcolumntype{L}[1]{>{\raggedright\arraybackslash}p{#1}}

% Senza il dizionario italiano la sillabazione automatica sbaglierebbe.
\hyphenpenalty=10000
\exhyphenpenalty=10000
\tolerance=3000
\emergencystretch=3em

\definecolor{scuro}{RGB}{25,50,90}
\definecolor{sfondo}{RGB}{246,247,250}
\definecolor{commento}{RGB}{110,120,135}
\definecolor{chiave}{RGB}{160,40,60}

\titleformat{\section}{\Large\bfseries\color{scuro}}{\thesection.}{0.6em}{}
\titleformat{\subsection}{\large\bfseries\color{scuro}}{\thesubsection}{0.6em}{}
\titlespacing{\section}{0pt}{1.5em}{0.5em}
\setlist[itemize]{topsep=3pt,itemsep=2pt,parsep=0pt,leftmargin=1.2em}

\lstset{
  language=C,
  basicstyle=\ttfamily\footnotesize,
  keywordstyle=\color{chiave}\bfseries,
  commentstyle=\color{commento}\itshape,
  backgroundcolor=\color{sfondo},
  frame=leftline,
  rulecolor=\color{scuro},
  framesep=6pt,
  xleftmargin=10pt,
  aboveskip=8pt,
  belowskip=8pt,
  showstringspaces=false,
  breaklines=true,
  columns=fullflexible,
}

\newcommand{\prima}[1]{\par\smallskip\noindent\textbf{\color{scuro}Prima:} #1\par}
\newcommand{\dopo}[1]{\par\smallskip\noindent\textbf{\color{scuro}Dopo:} #1\par\smallskip}
\newcommand{\cd}[1]{\texttt{\small #1}}
\newenvironment{blocco}{\par\noindent\begin{minipage}{\linewidth}}{\end{minipage}\par}
\newcommand{\etichetta}[1]{\par\smallskip\noindent{\small\bfseries\color{scuro}#1}\par\vspace{-4pt}}

\begin{document}

\begin{center}
{\LARGE\bfseries\color{scuro} Parallelizzazione con MPI del solutore\\[2pt]
Navier--Stokes--Brinkman}\\[8pt]
{\large Relazione sul lavoro svolto}\\[10pt]
\end{center}

\hrule
\vspace{1.2em}

\section{Di cosa si tratta}

Il programma risolve le equazioni di Navier--Stokes--Brinkman su una griglia
tridimensionale, con lo schema a separazione delle direzioni visto a lezione.
Finora girava su un solo processore.

L'obiettivo è farlo girare su più processori contemporaneamente usando MPI, la
libreria standard del calcolo parallelo. L'idea di fondo è semplice: dividere
la griglia in blocchi, assegnarne uno a ciascun processore, e far sì che quando
a un processore serve un dato che sta a casa del vicino, glielo chieda con un
messaggio.

Il lavoro è diviso in sette fasi, dalla A alla G. Le prime due preparano il
terreno senza cambiare il modo in cui il solutore calcola; le successive
collegano i pezzi, e alla fine il programma gira davvero su più processori e
si misura quanto ci si guadagna.

In questa relazione ogni pezzo è mostrato con il codice vero e spiegato subito
dopo.

\section{Perché non basta ``accendere'' MPI}

Ci sono due ostacoli, di natura molto diversa.

\subsection*{Primo ostacolo: il codice credeva di avere sempre tutta la griglia}

Le dimensioni della griglia erano tre costanti fissate al momento della
compilazione: \cd{WIDTH}, \cd{HEIGHT}, \cd{DEPTH}. Ogni ciclo, ogni calcolo di
posizione e ogni controllo del tipo ``sono sul bordo?'' le usava direttamente.

Dividendo la griglia, ciascun processore ne possiede solo una fetta. Il codice
non aveva alcun modo di distinguere \emph{la griglia intera} da \emph{la mia
parte}: per lui erano la stessa cosa.

\subsection*{Secondo ostacolo: l'algoritmo di Thomas è una catena}

Il solutore risolve sistemi tridiagonali con l'algoritmo di Thomas. Il valore in
un punto si calcola da quello del punto precedente, che si calcola dal
precedente ancora: è una catena rigidamente sequenziale.

Se spezziamo una linea fra due processori, il secondo non può cominciare finché
il primo non ha finito. Non avremmo parallelizzato niente: stesso tempo di
prima, più il costo dei messaggi. La soluzione, vista a lezione, è il
\emph{complemento di Schur}.

\section{Fase A --- La griglia diventa ``un pezzo''}

\subsection*{La struttura che sostituisce le costanti}

\prima{Tre costanti di compilazione, usate ovunque.}
\dopo{Una struttura dati, passata come parametro a ogni funzione che percorre
la griglia.}

\begin{lstlisting}
typedef struct Decomp {
    int n_global[3];   /* dimensione della griglia intera        */
    int n[3];          /* celle che possiedo io                  */
    int start[3];      /* indice globale della mia prima cella   */
    size_t stride[3];  /* distanza in memoria fra celle vicine   */
    size_t base;       /* offset della prima cella posseduta     */
    int is_first[3];   /* tocco la parete inferiore?             */
    int is_last[3];    /* tocco la parete superiore?             */
    size_t n_cells;    /* celle da allocare per un campo         */
} Decomp;
\end{lstlisting}

Con una griglia $32\times16\times24$ su un solo processore i campi valgono:
\cd{n\_global} e \cd{n} entrambi $\{32,16,24\}$, \cd{start} $\{0,0,0\}$,
\cd{stride} $\{1,32,512\}$, \cd{base} $0$, tutti i flag a $1$, e
\cd{n\_cells} $=12\,288$.

Il campo che porta tutto il peso è \cd{stride}. Significa: per spostarsi di una
cella lungo $Y$ bisogna avanzare di 32 posizioni nell'array, e per cambiare
piano di 512 (cioè $32\times16$). Averlo memorizzato come \emph{dato} invece di
ricostruirlo ogni volta dalle dimensioni è ciò che renderà indolore la fase C.

\subsection*{Come si raggiunge una cella}

\begin{lstlisting}
static inline size_t decomp_index(const Decomp *d, int i, int j, int k) {
    return d->base + (size_t)i * d->stride[0] + (size_t)j * d->stride[1] +
           (size_t)k * d->stride[2];
}

static inline int decomp_global(const Decomp *d, int i, int component) {
    return d->start[component] + i;
}
\end{lstlisting}

La prima funzione traduce una posizione nella griglia in una posizione in
memoria. Con i valori dell'esempio dà \cd{0 + i + 32*j + 512*k}, cioè
esattamente il calcolo che il codice faceva prima a mano: stessa aritmetica, ma
con i numeri presi dalla struttura.

La seconda traduce un indice \emph{locale} (``la mia terza cella'') in un indice
\emph{globale} (``la cella numero 15 del dominio''). Serve ogni volta che si
calcola una coordinata fisica o si chiede se siamo su una parete.

\subsection*{Come si riempie la struttura}

\begin{lstlisting}
static void decomp_set_layout(Decomp *d, int halo) {
    size_t allocated_x = (size_t)d->n[0] + 2 * (size_t)halo;
    size_t allocated_y = (size_t)d->n[1] + 2 * (size_t)halo;
    size_t allocated_z = (size_t)d->n[2] + 2 * (size_t)halo;

    d->stride[0] = 1;
    d->stride[1] = allocated_x;
    d->stride[2] = allocated_x * allocated_y;
    d->base = (size_t)halo * (d->stride[0] + d->stride[1] + d->stride[2]);
    d->n_cells = allocated_x * allocated_y * allocated_z;
}
\end{lstlisting}

Questa è l'unica funzione che decide dove stanno fisicamente i dati.
Il parametro \cd{halo} è il margine di celle lasciato su ogni faccia: oggi vale
zero, in fase C diventerà uno.

Il punto importante: \textbf{aggiungere il margine cambia solo questa
funzione}. Con \cd{halo} $=1$ gli stride diventano $\{1,34,612\}$ e \cd{base}
diventa $647$ invece di $0$ -- e le centinaia di chiamate a \cd{decomp\_index}
sparse nel codice continuano a funzionare senza essere toccate.

\subsection*{I quattro tipi di modifica}

Ogni singola modifica della fase A rientra in uno di questi quattro casi. I
primi tre compaiono tutti insieme in questo ciclo, che riempie un campo di
velocità.

\etichetta{Prima}
\begin{lstlisting}
size_t off = 0;
for (int k = 0; k < DEPTH; k++) {
  for (int j = 0; j < HEIGHT; j++) {
    for (int i = 0; i < WIDTH; i++) {
      Real x = centered_physical_coord(i, 0);
      ...
      v_x[off] = vector_fn(staggered_physical_coord(i, 0), y, z, time, 0);
      off++;
    }
  }
}
\end{lstlisting}

\etichetta{Dopo}
\begin{lstlisting}
for (int k = 0; k < d->n[2]; k++) {
  int gk = decomp_global(d, k, 2);
  for (int j = 0; j < d->n[1]; j++) {
    int gj = decomp_global(d, j, 1);
    size_t row = decomp_index(d, 0, j, k);
    for (int i = 0; i < d->n[0]; i++) {
      int gi = decomp_global(d, i, 0);
      Real x = centered_physical_coord(gi, 0);
      ...
      size_t off = row + (size_t)i;
      v_x[off] = vector_fn(staggered_physical_coord(gi, 0), y, z, time, 0);
    }
  }
}
\end{lstlisting}

Tre cose sono cambiate.

\begin{itemize}
\item \textbf{Gli estremi dei cicli.} Da ``scorro tutta la griglia'' a ``scorro
le celle che possiedo''.
\item \textbf{La posizione in memoria.} Il contatore \cd{off++} avanzava di una
posizione alla volta: funziona solo se le celle sono perfettamente consecutive.
Con il margine non lo saranno più -- dopo l'ultima cella di una riga vengono le
celle di contorno, poi la riga successiva -- e \cd{off++} finirebbe dentro il
margine. Ora la posizione si calcola.
\item \textbf{La coordinata fisica.} Da \cd{i} a \cd{decomp\_global(d, i, 0)}.
La posizione di una cella nel mondo fisico dipende da dove sta nel
\emph{dominio}, non da dove sta nel mio pezzo di memoria. Senza questa
conversione, il secondo processore calcolerebbe le forze e le condizioni al
contorno come se si trovasse all'origine del dominio.
\end{itemize}

Il quarto caso è il più insidioso, e riguarda i controlli sul bordo.

\etichetta{Prima}
\begin{lstlisting}
if (i < 1 || i >= WIDTH - 1 || j < 1 || j >= HEIGHT || ...) {
    return 0.0;
}
\end{lstlisting}

\etichetta{Dopo}
\begin{lstlisting}
int gi = decomp_global(d, i, 0);
int gj = decomp_global(d, j, 1);
if (gi < 1 || gi >= d->n_global[0] - 1 || gj < 1 || ...) {
    return 0.0;
}
\end{lstlisting}

Prima, ``ultima cella del mio array'' e ``sono sulla parete fisica'' erano la
stessa condizione. Con più processori le due cose si separano: un processore in
mezzo al dominio ha una sua ultima cella, ma lì non c'è nessuna parete, c'è il
vicino.

Sbagliare questo controllo crea \emph{pareti fantasma} dentro il fluido: la
simulazione gira, i numeri sembrano plausibili, e sono sbagliati. È il tipo di
errore che si scopre solo confrontando i risultati fra numeri di processori
diversi.

\section{Fase B --- I mattoni del parallelo}

La fase B è divisa in quattro passi. Due dei quattro non usano MPI affatto:
è voluto, e serve a separare la domanda ``la matematica è giusta?'' dalla
domanda ``la comunicazione è giusta?''. Alla fine della fase B tutti i pezzi
esistono e sono verificati singolarmente, ma il solutore non è ancora
parallelo: nessuno li ha ancora collegati.

\subsection*{B0 --- Un muro fra MPI e il solutore}

\prima{Nessun supporto per MPI.}
\dopo{Due file nuovi che sono l'unico posto del progetto a conoscere MPI.}

\begin{lstlisting}
/* parallel.h */
void par_init(int *argc, char ***argv);
void par_finalize(void);
int  par_rank(void);   /* chi sono io: da 0 a par_size()-1 */
int  par_size(void);   /* quanti siamo */
\end{lstlisting}

\begin{lstlisting}
/* parallel.c */
#ifdef USE_MPI
  #include <mpi.h>
  int par_rank(void) {
      int rank;
      MPI_Comm_rank(MPI_COMM_WORLD, &rank);
      return rank;
  }
#else
  int par_rank(void) { return 0; }   /* siamo un processo solo */
#endif
\end{lstlisting}

Il solutore chiama \cd{par\_rank()} e non sa che sotto ci sia
\cd{MPI\_Comm\_rank}. Non lo sa e non deve saperlo.

Il vantaggio non è estetico. Senza l'opzione di compilazione, le funzioni
diventano banali e il programma seriale continua a funzionare come prima, senza
MPI installato. La versione seriale resta così un \textbf{termine di paragone}:
ogni volta che il parallelo darà un numero sospetto, lo confronteremo con la
versione che sappiamo essere giusta.

\subsection*{B1 --- Spezzare la catena, senza ancora usare MPI}

Prima il metodo sequenziale, che resta il riferimento.

\begin{lstlisting}
void thomas_solve(int n, const Real *a, const Real *b, const Real *c,
                  const Real *f, Real *x, Real *scratch) {
    Real denominator = b[0];
    scratch[0] = c[0] / denominator;
    x[0] = f[0] / denominator;

    for (int i = 1; i < n; i++) {                    /* eliminazione */
        denominator = b[i] - a[i] * scratch[i - 1];
        scratch[i] = c[i] / denominator;
        x[i] = (f[i] - a[i] * x[i - 1]) / denominator;
    }

    for (int i = n - 2; i >= 0; i--) {               /* sostituzione */
        x[i] -= scratch[i] * x[i + 1];
    }
}
\end{lstlisting}

Si vede bene la catena: nel primo ciclo il valore in \cd{i} dipende da quello in
\cd{i-1}. Non c'è modo di calcolarli in parallelo.

\medskip
Il complemento di Schur rompe la catena così: dentro il mio blocco, la soluzione
è la somma di due contributi. Il primo è la risposta al carico che ho io, e lo
calcolo subito da solo. Il secondo è la risposta ai valori alle mie due
estremità, che ancora non conosco -- ma posso calcolare quanto ``pesa'' ciascuna
estremità sul mio interno, con due sole risoluzioni locali in più.

\begin{lstlisting}
/* ---- 1. ogni blocco per conto suo: tre risoluzioni locali ---- */
thomas_solve(len, a+begin, b+begin, c+begin, f+begin, y+begin, scratch);

if (p > 0) {                        /* influenza dell'estremita' sinistra */
    memset(rhs, 0, len * sizeof(Real));
    rhs[0] = a[begin];
    thomas_solve(len, a+begin, b+begin, c+begin, rhs, lft+begin, scratch);
}

if (has_right) {                    /* influenza dell'estremita' destra   */
    memset(rhs, 0, len * sizeof(Real));
    rhs[len - 1] = c[internal_end - 1];
    thomas_solve(len, a+begin, b+begin, c+begin, rhs, rgt+begin, scratch);
}
\end{lstlisting}

\cd{y} è la risposta al carico proprio; \cd{lft} e \cd{rgt} sono le due funzioni
di influenza, ottenute risolvendo lo stesso sistema con un termine noto che vale
zero ovunque tranne che a un'estremità. Nessuna di queste tre risoluzioni guarda
i dati di un altro blocco.

\medskip
A quel punto le uniche incognite rimaste sono i punti di giunzione fra blocchi.
Sostituendo l'espressione della soluzione interna nelle equazioni scritte in
quei punti si ottiene un sistema nelle sole giunzioni, che risulta essere di
nuovo tridiagonale:

\begin{lstlisting}
/* ---- 2. il sistema rimasto sulle giunzioni ---- */
int m = end - 1;         /* il punto di giunzione                   */
int left = m - 1;        /* ultimo punto interno del blocco p       */
int right = next_begin;  /* primo punto interno del blocco p+1      */

ra[p] = -a[m] * lft[left];
rb[p] = b[m] - a[m] * rgt[left] - c[m] * lft[right];
rc[p] = -c[m] * rgt[right];
rf[p] = f[m] - a[m] * y[left] - c[m] * y[right];

thomas_solve(interfaces, ra, rb, rc, rf, rx, scratch);
\end{lstlisting}

Con quattro blocchi questo sistema ha tre incognite, contro le centoventotto di
partenza: è minuscolo. Risolto quello, ogni blocco ricompone la sua parte:

\begin{lstlisting}
/* ---- 3. ogni blocco di nuovo per conto suo ---- */
Real x_left  = (p > 0)    ? rx[p - 1] : (Real)0;
Real x_right = has_right  ? rx[p]     : (Real)0;

for (int i = begin; i < internal_end; i++) {
    x[i] = y[i] - x_left * lft[i] - x_right * rgt[i];
}
\end{lstlisting}

Questa riga è la sovrapposizione degli effetti scritta in codice: la risposta al
mio carico, meno il contributo di ciascuna estremità pesato per il valore che
quell'estremità ha effettivamente assunto.

Il punto importante: \textbf{non è un'approssimazione}. È lo stesso identico
risultato del metodo sequenziale, con le operazioni fatte in un ordine diverso.

In questo passo i ``blocchi'' sono ancora fette di un unico array dentro un
unico processo, e non si comunica niente. Serve a verificare la matematica
quando è ancora ispezionabile con un debugger ordinario.

\subsection*{B2 --- Dare a ciascun processore la sua fetta}

\prima{I processori esistevano ma nessuno diceva loro quale parte della griglia
gli spettasse.}
\dopo{I processori sono disposti su una griglia tridimensionale, e ognuno ne
ricava la propria fetta.}

\begin{lstlisting}
void par_topology_init(const int procs[3]) {
    int dims[3] = {procs[0], procs[1], procs[2]};
    int periods[3] = {0, 0, 0};

    MPI_Dims_create(par_size(), 3, dims);
    MPI_Cart_create(MPI_COMM_WORLD, 3, dims, periods, 0, &cart_comm);
    MPI_Cart_get(cart_comm, 3, cart_dims, periods, cart_coords);

    for (int c = 0; c < 3; c++) {          /* un gruppo per direzione */
        int keep[3] = {0, 0, 0};
        keep[c] = 1;
        MPI_Cart_sub(cart_comm, keep, &line_comm[c]);
    }
}
\end{lstlisting}

\cd{MPI\_Dims\_create} riempie solo le direzioni lasciate a zero: passando
$\{1,1,0\}$ si ottiene una divisione a fette lungo $Z$, passando $\{0,0,0\}$ è
MPI a scegliere una forma equilibrata (con otto processori, $2\times2\times2$).

L'ultimo ciclo crea tre gruppi di processori, uno per direzione: sono i gruppi
che si scambieranno le giunzioni nel complemento di Schur.

\begin{lstlisting}
void decomp_init_mpi(Decomp *d) {
    int dims[3], coords[3];
    par_dims(dims);
    par_coords(coords);

    for (int c = 0; c < 3; c++) {
        int begin, end;
        decomp_share(d->n_global[c], dims[c], coords[c], &begin, &end);

        d->n[c] = end - begin;
        d->start[c] = begin;
        d->is_first[c] = (coords[c] == 0);
        d->is_last[c] = (coords[c] == dims[c] - 1);
    }
    decomp_set_layout(d, 0);
}
\end{lstlisting}

Qui si vede il senso della fase A: è \emph{l'unica} funzione che cambia rispetto
alla versione seriale. \cd{n} diventa più piccolo, \cd{start} non è più zero, e
i flag delle pareti diventano veri solo per i processori che stanno davvero agli
estremi. Tutto il resto del solutore non se ne accorge.

Con 24 celle lungo $Z$ e cinque processori i blocchi escono di dimensione
$5,5,5,5,4$: il resto della divisione viene distribuito una cella per volta ai
primi blocchi, così le dimensioni differiscono al massimo di uno.

\subsection*{B3 --- Il complemento di Schur su processori veri}

\prima{Il complemento di Schur funzionava su fette di un array, in un solo
processo.}
\dopo{Funziona con un blocco per processore, e i valori che prima si leggevano
dall'array del vicino ora viaggiano come messaggi.}

Le fasi 1 e 3 restano identiche a prima e non comunicano. Cambia solo la fase 2:

\begin{blocco}
\begin{lstlisting}
/* Il mio vicino di sinistra ha bisogno di sapere quanto vale
   il mio primo punto interno nelle tre risposte locali.        */
Real mine[3] = {y[0], lft[0], rgt[0]};
Real from_right[3] = {0, 0, 0};
par_shift_real(axis, -1, mine, from_right, 3);

/* La riga del sistema sulle giunzioni che spetta a me.
   L'ultimo processore non ha giunzioni e manda una riga innocua. */
Real row[4] = {0, 1, 0, 0};

if (has_right) {
    int m = n_local - 1;
    int left = len - 1;
    row[0] = -a[m] * lft[left];
    row[1] = b[m] - a[m] * rgt[left] - c[m] * from_right[1];
    row[2] = -c[m] * from_right[2];
    row[3] = f[m] - a[m] * y[left] - c[m] * from_right[0];
}

/* Tutti ricevono tutte le righe, e tutti risolvono il sistemino:
   cosi' nessuno deve rispedire indietro la soluzione.           */
par_line_allgather(axis, row, 4, rows);
\end{lstlisting}
\end{blocco}

Confrontando con il codice di B1 si vede che le formule sono le stesse: dove
prima c'era \cd{lft[right]} -- un valore letto dall'array del blocco vicino --
ora c'è \cd{from\_right[1]}, cioè lo stesso valore arrivato per posta.

Lo scambio col vicino è implementato così:

\begin{lstlisting}
void par_shift_real(int axis, int step,
                    const Real *send, Real *recv, int count) {
    MPI_Sendrecv(send, count, PAR_REAL, mpi_neighbor(axis,  step), 0,
                 recv, count, PAR_REAL, mpi_neighbor(axis, -step), 0,
                 cart_comm, MPI_STATUS_IGNORE);
}
\end{lstlisting}

Due dettagli meritano una nota.

\cd{MPI\_Sendrecv} invece di una \cd{Send} seguita da una \cd{Recv}: due
processori affacciati devono scambiarsi dati nello stesso istante, e con
chiamate separate si bloccherebbero a vicenda non appena i messaggi superano la
soglia oltre la quale MPI smette di conservarli in un'area temporanea.

\cd{mpi\_neighbor} restituisce un valore speciale, \cd{MPI\_PROC\_NULL}, dove il
vicino non c'è. MPI tratta un invio o una ricezione verso il vuoto come
un'operazione che non fa niente, quindi non serve nessun caso particolare per i
processori sul bordo.

\medskip
Il conto totale della comunicazione è sorprendentemente piccolo:

\begin{center}
\begin{tabular}{@{}L{7.2cm}L{7.2cm}@{}}
\toprule
\textbf{Fase} & \textbf{Comunicazione} \\
\midrule
1. Calcolo locale (tre risoluzioni) & nessuna \\
2. Sistema sulle giunzioni & 3 numeri col vicino, 4 per processore \\
3. Ricomposizione locale & nessuna \\
\bottomrule
\end{tabular}
\end{center}

\medskip
Tutto il resto è calcolo che ciascun processore fa per conto proprio: è
esattamente ciò che serve perché il parallelo convenga.

\section{Fase C --- Le celle in prestito dal vicino}

\prima{Ogni processo possedeva solo le proprie celle. Ma per calcolare vicino
al bordo del proprio blocco servono i dati di quello accanto.}

\dopo{Ogni processo alloca un anello di celle in più tutt'intorno, dove tiene
una copia dell'ultima fila del vicino, e se la fa aggiornare con un messaggio
prima di usarla.}

Il conto che il solutore fa in ogni cella guarda anche la cella prima e quella
dopo: è una derivata seconda. Per una cella sul bordo del blocco, ``quella
dopo'' appartiene al vicino. Le celle in prestito servono a questo, e si
chiamano \emph{celle fantasma}.

\subsection*{Allocarle è costato una riga}

\begin{lstlisting}
decomp_set_layout(d, 1);   /* prima era 0: nessun margine */
\end{lstlisting}

Questo è il ritorno dell'investimento della fase A. Il margine cambia dove
stanno i dati in memoria -- le righe non sono più contigue, la prima cella non
è più all'indice zero -- ma tutte le centinaia di calcoli di posizione sparsi
nel codice passano da \cd{decomp\_index}, che legge quei numeri dalla
struttura. Non è stato necessario toccarne nemmeno uno.

\subsection*{Lo scambio}

\begin{lstlisting}
for (int axis = 0; axis < 3; axis++) {
    int below = par_neighbor(axis, -1);
    int above = par_neighbor(axis, +1);

    if (below == PAR_NO_NEIGHBOR && above == PAR_NO_NEIGHBOR) {
        continue;   /* direzione non divisa: niente da scambiare */
    }

    /* La mia ultima faccia va a chi sta sopra; dal basso arriva la sua. */
    face_copy(d, field, axis, d->n[axis] - 1, packet, 1);
    par_shift_real(axis, +1, packet, arrived, face);
    if (below != PAR_NO_NEIGHBOR) {
        face_copy(d, field, axis, -1, arrived, 0);
    }

    /* E specularmente verso il basso. */
    face_copy(d, field, axis, 0, packet, 1);
    par_shift_real(axis, -1, packet, arrived, face);
    if (above != PAR_NO_NEIGHBOR) {
        face_copy(d, field, axis, d->n[axis], arrived, 0);
    }
}
\end{lstlisting}

\cd{face\_copy} copia una faccia del blocco dentro un pacchetto contiguo, o
viceversa: impacchetta e scarica. \cd{par\_shift\_real} è la stessa funzione
già usata nella fase B, quindi non c'è nessuna primitiva MPI nuova.

Due dettagli che vale la pena spiegare.

MPI mette a disposizione i \emph{tipi derivati} proprio per descrivere dati non
contigui senza copiarli a mano. Non li ho usati perché, come dice la
documentazione stessa, in mancanza di hardware apposito la libreria fa
comunque questa identica copia in un'area contigua. Farla a mano costa uguale
e si legge; in più vale subito per tutte e tre le direzioni.

Non servono gli \emph{angoli}. I conti guardano una cella più in là in
\emph{una} direzione alla volta, mai in diagonale: bastano le sei facce, non
serve un secondo giro per gli spigoli.

\subsection*{Quanto si scambia}

Su una griglia $128^3$ divisa fra quattro processi, la faccia da spedire è
128\,KB contro i 4\,MB del blocco locale: si scambia circa il \textbf{3\%} dei
dati per poter calcolare il restante 97\% in autonomia. È il rapporto fra
superficie e volume, ed è la ragione per cui dividere il dominio funziona.

\section{Fase D --- Il primo calcolo davvero parallelo}

\prima{Tutti i pezzi esistevano, ma nessuno li aveva collegati: il solutore
girava ancora come se fosse solo.}

\dopo{I sistemi lungo Z passano dal complemento di Schur, le celle di contorno
vengono aggiornate al momento giusto, e il risultato non dipende più da quanti
processi si usano.}

\subsection*{Scrivere il sistema invece di risolverlo di corsa}

Prima il codice costruiva e risolveva il sistema tridiagonale in un unico
ciclo, senza mai scrivere da nessuna parte le tre diagonali. Per consegnarlo
al complemento di Schur bisogna invece scriverle:

\begin{lstlisting}
if (along == 0) {
    /* Parete inferiore del dominio: il valore e' imposto. */
    lower[at] = 0.0;
    diagonal[at] = 1.0;
    upper[at] = 0.0;
    known[at] = bc_left(data->bc_velocity, gi, gj, gk, t_step, v_comp);
} else if (along < last_global) {
    lower[at] = w_i;
    diagonal[at] = 1.0 - 2.0 * w_i;      /* riga interna */
    upper[at] = w_i;
    known[at] = rhs;
} else {
    ...                                   /* parete superiore */
}
\end{lstlisting}

La variabile chiave è \cd{along}: è l'indice \textbf{globale}, cioè la
posizione della cella nel dominio intero. Un processo in mezzo non ha nessuna
parete, quindi per lui \cd{along} non vale mai zero e usa ovunque la riga
interna.

Questo è il punto in cui il quarto tipo di modifica della fase A -- quello
delle ``pareti fantasma'' -- viene finalmente messo alla prova. Se il test
fosse rimasto sull'indice locale, ogni processo avrebbe creduto di avere una
parete al proprio bordo.

\subsection*{Il momento in cui aggiornare le celle di contorno}

\begin{lstlisting}
for (int t_step = 1; t_step <= STEPS; t_step++) {
    refresh_vector_halo(decomp, &solver_mem_state->eta);
    refresh_vector_halo(decomp, &solver_mem_state->zeta);
    refresh_vector_halo(decomp, &solver_mem_state->u);
    par_exchange_halo(decomp, solver_mem_state->pressure_star.v);

    momentum_step(...);

    /* La velocita' e' appena cambiata e il passo di pressione la rilegge. */
    refresh_vector_halo(decomp, &solver_mem_state->u);
    pressure_step(...);
}
\end{lstlisting}

Due soli momenti per passo temporale: all'inizio, e in mezzo dopo che la
velocità è stata aggiornata. La regola è semplice: si aggiorna un campo quando
è cambiato e sta per essere riletto oltre il proprio bordo.

\subsection*{Una comunicazione per gruppo di linee, non per linea}

Questa è la scelta di progetto più importante di tutta la fase.

Il solutore chiama il complemento di Schur \emph{una volta per ogni linea della
griglia}. Su $128^3$ sono circa sedicimila linee. Un messaggio per linea
significherebbe sedicimila messaggi per ogni risoluzione: costerebbe molto più
del calcolo che permette di fare.

La soluzione è fare prima tutto il lavoro locale di un intero gruppo di linee,
e poi scambiare una volta sola per l'intero gruppo:

\begin{lstlisting}
/* 1. ogni linea per conto suo, nessuna comunicazione */
for (int l = 0; l < lines; l++) {
    thomas_solve(len, a + off, b + off, c + off, f + off, y + off, scratch);
    ...   /* piu' le due funzioni di influenza */
}

/* 2. una sola volta, per tutte le linee insieme */
par_shift_real(axis, -1, mine, from_right, 3 * lines);
par_line_allgather(axis, row, 4 * lines, rows);

/* 3. ogni linea di nuovo per conto suo */
\end{lstlisting}

Si passa da migliaia di messaggi minuscoli a poche centinaia per passo
temporale.

\subsection*{Il controllo che chiude la fase}

Il complemento di Schur è un metodo \emph{esatto}: non approssima niente,
riordina soltanto le stesse operazioni. Quindi il numero di processi non deve
cambiare il risultato. Il confronto con la versione a un processo solo:

\begin{center}
\begin{tabular}{@{}lll@{}}
\toprule
& \textbf{un processo} & \textbf{da 1 a 8 processi} \\
\midrule
errore su $u_x$ & 5.8196668537e-04 & identico, cifra per cifra \\
errore su $u_y$ & 5.4026976672e-04 & identico, cifra per cifra \\
errore su $u_z$ & 2.8691858457e-04 & identico, cifra per cifra \\
errore su $p$   & 1.9758785661e-01 & identico, cifra per cifra \\
\bottomrule
\end{tabular}
\end{center}

\medskip
Non era scontato ottenere un'identità \emph{esatta}: cambiando l'ordine delle
somme in virgola mobile ci si aspetterebbero scarti nelle ultime cifre. Il
fatto che non ce ne siano conferma che il metodo distribuito fa anche le stesse
operazioni nello stesso ordine.

\subsection*{Una cosa anticipata dalla fase F}

Al primo tentativo gli errori sembravano \emph{migliorare} aumentando i
processi. Non era vero: ogni processo calcolava la norma dell'errore solo sul
proprio blocco, senza sommare quella degli altri. È bastato aggiungere tre
somme globali, ma senza di esse la fase D non aveva un modo per verificarsi.

\section{Fase E --- Dividere anche nelle altre due direzioni}

\prima{La griglia si divideva solo lungo Z: le altre due direzioni davano
risultati sbagliati, e il programma si fermava con un messaggio se qualcuno ci
provava.}

\dopo{Si divide in tutte e tre le direzioni, e la forma della griglia di
processi la sceglie MPI.}

\subsection*{La scoperta che ha guidato il lavoro}

Servivano le versioni con Schur anche per X e Y. Ma i tre passi sono la
\textbf{stessa operazione ripetuta su assi diversi} --- ed è esattamente come
li presenta il corso:

$$(I-\gamma\partial_{xx})\Delta\eta = \xi-\eta^n, \qquad
  (I-\gamma\partial_{yy})\Delta\zeta = \eta^{n+1}-\zeta^n, \qquad
  (I-\gamma\partial_{zz})\Delta u = \zeta^{n+1}-u^n$$

Scriverne altre due copie avrebbe voluto dire tre funzioni quasi identiche. Ne
ho scritta \textbf{una sola}, che prende l'asse come parametro:

\begin{lstlisting}
/* Le linee corrono lungo `axis`; delle altre due direzioni, `group`
 * raccoglie le linee risolte insieme, `outer` e' il ciclo esterno. */
const int group = (axis == 0) ? 1 : 0;
const int outer = (axis == 2) ? 1 : 2;
const int length = d->n[axis];
const int lines  = d->n[group];
const size_t step = d->stride[axis];
\end{lstlisting}

\cd{step} è la distanza in memoria fra due celle vicine lungo quell'asse: 1
lungo X, la larghezza di una riga lungo Y, quella di un piano lungo Z. Prendendola
dalla struttura, lo stesso ciclo percorre linee in qualunque direzione.

\subsection*{Lo stesso per la pressione}

Anche i tre passi della pressione hanno la \emph{stessa matrice}: cambiano solo
il campo da cui si legge e quello in cui si scrive. Una funzione sola, e la
cascata del corso diventa leggibile in cinque righe:

\begin{lstlisting}
compute_div(decomp, buffer, &solver_mem_state->u);
pressure_direction(decomp, 0, buffer, star);   /* psi */
pressure_direction(decomp, 1, star, buffer);   /* phi */
pressure_direction(decomp, 2, buffer, star);   /* fi  */
update_pressure(decomp, solver_mem_state);
\end{lstlisting}

I due campi si scambiano il ruolo a ogni passo, quindi ne bastano due.

\subsection*{Il risultato}

La fase E ha \textbf{tolto} codice invece di aggiungerne:

\begin{center}
\begin{tabular}{@{}L{5cm}L{4.4cm}L{4.4cm}@{}}
\toprule
& \textbf{Prima} & \textbf{Dopo} \\
\midrule
Quantità di moto & 3 funzioni, 297 righe & 1 funzione, 150 righe \\
Pressione & 4 funzioni, 250 righe & 1 funzione, 140 righe \\
\bottomrule
\end{tabular}
\end{center}

\medskip
E non è servito toccare nulla in \cd{parallel.c}: lo scambio delle celle di
contorno era già scritto per tutti e tre gli assi fin dalla fase C, e il
complemento di Schur prende l'asse come parametro fin dalla fase B. Era il
motivo di quelle scelte.

\section{Fase F --- Scrivere i risultati in parallelo}

\prima{Ogni processo avrebbe scritto il proprio pezzo di dominio in un file con
lo stesso nome: si sarebbero sovrascritti a vicenda.}

\dopo{Ogni processo scrive il proprio file, e il primo scrive un indice che
tiene insieme i pezzi.}

L'indice è un file di poche righe che elenca dove sta ciascun pezzo. Aprendo
quello in ParaView si vede il dominio intero ricomposto:

\begin{lstlisting}
<PImageData WholeExtent="0 23 0 15 0 11" GhostLevel="0" ...>
  <Piece Extent="0 7 0 7 0 11"    Source="sol_0005_p000.vti"/>
  <Piece Extent="0 7 8 15 0 11"   Source="sol_0005_p001.vti"/>
  <Piece Extent="8 15 0 7 0 11"   Source="sol_0005_p002.vti"/>
  ...
\end{lstlisting}

\subsection*{Nessuno raccoglie i dati}

Non c'è nessuna raccolta su un processo solo: sarebbe il collo di bottiglia
classico, e su griglie grandi significherebbe far transitare gigabyte
attraverso un unico processo. Ognuno scrive per conto proprio, in parallelo.

E il processo che scrive l'indice \textbf{non chiede a nessuno dove stanno gli
altri}:

\begin{lstlisting}
for (int rank = 0; rank < par_size(); rank++) {
    int coords[3], begin[3], end[3];

    par_coords_of(rank, coords);
    for (int c = 0; c < 3; c++) {
        decomp_share(d->n_global[c], dims[c], coords[c], &begin[c], &end[c]);
    }
    fprintf(fp, "    <Piece Extent=\"%d %d %d %d %d %d\" Source=\"%s\"/>\n", ...);
}
\end{lstlisting}

Conoscendo la forma della griglia di processi, ricalcola i blocchi altrui con
\cd{decomp\_share}: \textbf{la stessa funzione con cui ciascuno si è preso il
suo}. Una sola regola usata in due posti, quindi l'indice e i pezzi non possono
contraddirsi.

\section{Fase G --- Quanto conviene}

Le misure sono state fatte su un portatile con processore AMD Ryzen~7 7435HS:
\textbf{8 nuclei fisici} (16 con il multithreading) e 7\,GB di memoria. Ogni
configurazione è stata ripetuta tre volte tenendo il tempo migliore: senza
ripetizioni il rumore di fondo della macchina rendeva i numeri inaffidabili.

\subsection*{Le due domande}

Lo \emph{strong scaling} chiede: a parità di problema, quanto si accorcia il
tempo aggiungendo processori? Il \emph{weak scaling} chiede: a parità di lavoro
per processore, il tempo resta lo stesso mentre il problema cresce?

\begin{center}
\begin{tabular}{@{}rrrr@{\hskip 2.2em}rrr@{}}
\toprule
& \multicolumn{3}{c}{\textbf{Strong} (problema fisso $128^3$)}
& \multicolumn{3}{c}{\textbf{Weak} ($64^3$ per processo)} \\
\cmidrule(r){2-4}\cmidrule(l){5-7}
\textbf{proc.} & tempo/passo & speedup & effic. & tempo/passo & rapporto & effic. \\
\midrule
1 & 941 ms & 1.00 & 100\% & 123 ms & 1.00 & 100\% \\
2 & 551 ms & 1.71 &  85\% & 149 ms & 0.83 &  83\% \\
4 & 340 ms & 2.77 &  69\% & 166 ms & 0.74 &  74\% \\
8 & 228 ms & \textbf{4.13} & 52\% & 232 ms & 0.53 & \textbf{53\%} \\
\bottomrule
\end{tabular}
\end{center}

\subsection*{Da dove viene la perdita: tre misure, non tre ipotesi}

Invece di attribuire la perdita a un generico ``sovraccarico'', l'ho scomposta.

\textbf{Contesa di memoria.} Ho lanciato \textbf{otto copie indipendenti} dello
stesso problema, senza \emph{nessuna} comunicazione fra loro. Rallentamento:
\textbf{1{,}35 volte}. È il limite della macchina, non del programma: otto
nuclei che si contendono un unico canale verso la memoria. Per un codice come
questo, che sposta molti più dati di quanti conti faccia, è il tetto vero.

\textbf{Il tempo dentro MPI.} Misurato direttamente, con un cronometro attorno
alle chiamate: al massimo carico sono 37\,ms su 232, cioè circa il
\textbf{16\%}.

\textbf{La vettorizzazione che si perde.} Questa è la scoperta più
interessante. Confrontando i singoli stadi a parità di lavoro per processo:

\begin{center}
\begin{tabular}{@{}lrrr@{}}
\toprule
\textbf{stadio} & \textbf{1 processo} & \textbf{8 processi} & \\
\midrule
passo lungo X & 61{,}0 ms & 101{,}9 ms & $\times$1{,}7 \\
passo lungo Y &  3{,}7 ms &  33{,}0 ms & $\times$8{,}9 \\
passo lungo Z &  3{,}9 ms &  37{,}9 ms & $\times$9{,}6 \\
\bottomrule
\end{tabular}
\end{center}

\medskip
Quasi dieci volte, e non è colpa né di MPI né della memoria. Il progetto
conteneva già dei nuclei di calcolo \emph{vettorizzati}, che risolvono una
linea intera in un colpo solo e sono molto veloci. Ma funzionano solo su una
linea intera: appena una direzione viene divisa fra processi, quella scorciatoia
non è più utilizzabile e si ricade sulla versione ordinaria.

Con otto processi disposti $2\times2\times2$ tutte e tre le direzioni sono
divise, quindi la vettorizzazione non si usa mai. Rifacendo lo studio con la
vettorizzazione spenta ovunque -- confronto alla pari -- lo speedup passa da
3{,}18 a \textbf{4{,}13}.

\medskip
Il bilancio a otto processi: circa 35\% contesa di memoria, 16\% comunicazione,
il resto il costo intrinseco del metodo di Schur, che fa tre risoluzioni locali
al posto di una.

\subsection*{I grafici}

Le curve sono in \cd{docs/scaling/scaling.svg}, prodotte da
\cd{scripts/plot\_scaling.py} a partire dai dati di
\cd{scripts/run\_scaling.sh}. Tre pannelli: lo speedup a confronto con la retta
ideale, l'efficienza del weak scaling, e le barre che mostrano quanta parte del
tempo se ne va dentro MPI.

\section{Come abbiamo verificato}

Ogni passo ha il suo controllo automatico, che risponde ``passato'' o
``fallito'' invece di stampare numeri da interpretare a occhio.

\begin{center}
\begin{tabular}{@{}L{3.5cm}L{7.2cm}L{3.5cm}@{}}
\toprule
\textbf{Controllo} & \textbf{Che cosa verifica} & \textbf{Esito} \\
\midrule
Risultati invariati &
Il solutore dà gli stessi numeri di prima della fase A &
identici, cifra per cifra \\[4pt]
Disposizione in memoria &
Cambiando come i dati stanno in RAM, il risultato non cambia &
zero celle diverse \\[4pt]
Schur (locale) &
Stessa risposta del metodo sequenziale, 24 configurazioni &
scarto $\sim 10^{-17}$ \\[4pt]
Fette dei processori &
I blocchi coprono la griglia una volta sola &
zero buchi, zero sovrapposizioni \\[4pt]
Schur (distribuito) &
Stessa risposta del sequenziale, con i messaggi &
scarto $\sim 10^{-17}$ \\[4pt]
Celle di contorno &
L'anello contiene i dati del vicino, e resta intatto dove c'è la parete &
zero celle sbagliate \\[4pt]
\textbf{Numero di processi} &
\textbf{Stessi errori di soluzione con 1, 2, 3, 4, 6, 8 processi e 8 forme di divisione} &
\textbf{identici, cifra per cifra} \\[4pt]
File di uscita &
I pezzi elencati nell'indice coprono il dominio una volta sola &
zero buchi, zero sovrapposizioni \\
\bottomrule
\end{tabular}
\end{center}

\medskip
Il settimo è quello decisivo, ed è stato ripetuto su due problemi di prova, due
griglie, con e senza vettorizzazione, e con otto forme diverse di divisione fra
cui i casi che dividono \emph{solo} lungo X e \emph{solo} lungo Y.

\subsection*{Abbiamo verificato anche i controlli stessi}

Un controllo che passa sempre non dimostra niente. Per esserne certi abbiamo
introdotto di proposito degli errori nel codice, uno alla volta, e verificato
che venissero scoperti. Su \textbf{trentuno} errori introdotti, trenta sono
stati scoperti; il trentunesimo era un \emph{mutante equivalente}, cioè un
programma scritto diversamente che calcola però la stessa identica cosa.

Questa verifica ha ripagato subito, perché ha rivelato un difetto \emph{nei
controlli stessi}. Il calcolo dello scarto massimo era scritto così:

\begin{lstlisting}
if (difference > worst) {
    worst = difference;
}
\end{lstlisting}

Rompendo di proposito una comunicazione, il risultato diventava ``non
numerico'' -- il valore che si ottiene, ad esempio, dividendo per zero. Ma per
quei valori \emph{ogni confronto risulta falso}: il valore non entrava mai nel
massimo, e il controllo dichiarava ``passato'' su un risultato completamente
rotto. La correzione è invertire il confronto:

\begin{lstlisting}
if (!(difference <= worst)) {
    worst = difference;
}
\end{lstlisting}

Sui numeri normali le due scritture sono equivalenti; su un valore non numerico
la seconda risulta vera, quindi l'errore emerge.

\section{Costi e limiti, dichiarati}

\textbf{Il codice seriale è diventato più lento del 12\%.} Per poter consegnare
il sistema al complemento di Schur bisogna scrivere le tre diagonali in
memoria, mentre prima venivano calcolate al volo. È il prezzo della
parallelizzazione, ed è misurato: da 37{,}5 a 42{,}2 unità di tempo per
cella-passo.

\textbf{La vettorizzazione e la divisione oggi competono.} Ogni direzione
divisa è una direzione su cui la scorciatoia vettorizzata non si può più usare.
Non è un difetto nascosto: è un lavoro non ancora fatto, e la strada è nota
(vettorizzare il percorso di Schur, le cui linee sono già indipendenti fra
loro).

\textbf{Il complemento di Schur fa tre risoluzioni locali invece di una}: la
risposta al carico proprio più le due funzioni di influenza. Per la pressione,
però, la matrice è costante nel tempo e uguale per tutte le linee: le due
funzioni di influenza si potrebbero calcolare una volta sola all'avvio. È
esattamente il ``preprocessing'' descritto a lezione, ed è l'ottimizzazione
successiva più ovvia.

\textbf{Un portatile non è la macchina giusta per uno studio di scaling.} Con
un solo canale verso la memoria, già a otto processi il limite è la banda e non
il programma. Gli stessi script funzionano senza modifiche su un calcolatore
con più nodi, dove i numeri sarebbero considerevolmente migliori.

\section{Cosa si potrebbe fare dopo}

\begin{description}[leftmargin=2.2em,style=nextline,itemsep=5pt]
\item[Vettorizzare il percorso parallelo] È il guadagno più grande a portata di
mano: recupererebbe il fattore dieci perso sui passi lungo Y e Z quando quelle
direzioni sono divise.

\item[Precalcolare le funzioni di influenza della pressione] La matrice non
cambia mai: le due risoluzioni in più si farebbero una volta all'avvio invece
che a ogni passo.

\item[Sovrapporre comunicazione e calcolo] MPI permette di avviare un messaggio
e continuare a calcolare mentre viaggia. Va misurato però se conviene davvero:
la documentazione stessa avverte che senza una scheda di rete che se ne occupi
da sola il guadagno spesso non c'è.

\item[Misurare su un calcolatore vero] Gli script sono pronti; servirebbe solo
una macchina con più nodi per vedere fin dove arriva il metodo.
\end{description}

\vspace{1em}
\hrule
\vspace{0.6em}
\noindent{\small Il solutore gira su più processori, la griglia si divide nelle
tre direzioni, e il risultato non dipende da quanti processori si usano: la
stessa soluzione, cifra per cifra, da uno a otto. Ogni pezzo è stato costruito
e verificato singolarmente prima di essere collegato agli altri, e quando
qualcosa non funzionava sapevamo dove non guardare.}

\end{document}
